\unit{Exposé 149 — Duality Theorems for Coherent Algebraic Sheaves}
\sid{149}
\src{169}
\seminaren{(May 1957)}
\paperhead{DUALITY THEOREMS FOR COHERENT ALGEBRAIC SHEAVES}%
  {by Alexandre GROTHENDIECK}

The results that follow, inspired by Serre's \fquote{algebraic duality theorem},
were found in the winter of 1955 and the winter of 1956. They are established
very simply with the aid of fairly elementary results on the cohomology of
projective spaces [3], and through intensive use of the homological algebra of
Cartan--Eilenberg, in the form given in [2].

\fsection{1}{The Ext's of sheaves of modules ({[2]}, Chapters 3 and 4)}

Let $X$ be a topological space equipped with a sheaf $\Osh$ of rings with
identity (not necessarily commutative). We consider the abelian category
$\Ccat^{\Osh}$ of sheaves of $\Osh$-modules, also called $\Osh$-Modules. Every
object of this category is known to admit an injective resolution, which makes
it possible to define in it the $\operatorname{Ext}$ functors having the familiar
formal properties. More precisely, to avoid confusion, we denote by
$\operatorname{Hom}_{\Osh}(X;A,B)$, or simply $\operatorname{Hom}(X;A,B)$, the
abelian group of $\Osh$-homomorphisms from $A$ to $B$, whereas
$\Homsh_{\Osh}(A,B)$ will denote the sheaf of germs of homomorphisms from $A$ to
$B$ ($A,B\in\Ccat^{\Osh}$). For fixed $A\in\Ccat^{\Osh}$, we define functors
$h_A$ and $\hsh_A$, taking values respectively in the category $\Ccat$ of
abelian groups and the category $\Ccat^{\Zsh}$ of sheaves of abelian groups on
$X$, by the formulas
\begin{equation}\tag{1.1}
  h_A(B)=\operatorname{Hom}_{\Osh}(X;A,B),
  \qquad
  \hsh_A(B)=\Homsh_{\Osh}(A,B).
\end{equation}

The functors $h_A$ and $\hsh_A$ are covariant left-exact functors. We consider
their right derived functors, denoted respectively by
$\operatorname{Ext}_{\Osh}^{p}(X;A,B)$ and $\Extsh_{\Osh}^{p}(A,B)$. Thus, by
definition,
\begin{equation}\tag{1.2}
\left\{
\begin{aligned}
\operatorname{Ext}_{\Osh}^{p}(X;A,B)
  &= (R^{p}h_A)(B)
   = H^{p}\bigl(\operatorname{Hom}_{\Osh}(X;A,C(B))\bigr),\\
\Extsh_{\Osh}^{p}(A,B)
  &= (R^{p}\hsh_A)(B)
   = H^{p}\bigl(\Homsh_{\Osh}(A,C(B))\bigr),
\end{aligned}
\right.
\end{equation}

where the symbol $R^{p}$ denotes passage to right derived functors, and where
$C(B)$ denotes an arbitrary injective resolution of $B$ in $\Ccat^{\Osh}$. Let
\[
  \Gamma:\Ccat^{\Zsh}\longrightarrow\Ccat
\]
denote the \fquote{sections} functor. Recall that its right derived functors are
denoted by
\[
  B\longmapsto H^{p}(X,B).
\]

